\newcommand{\DoNotLoadEpstopdf}{}
\documentclass[11pt,letterpaper]{article}
\usepackage[T1]{fontenc}
\usepackage{lmodern}
\usepackage[margin=1in]{geometry}
\usepackage{amsmath,amssymb,amsthm,mathtools}
\usepackage{enumitem}
\usepackage{microtype}
\usepackage{xurl}
\usepackage[colorlinks=true,linkcolor=blue,citecolor=blue,urlcolor=blue]{hyperref}
\usepackage{tikz-cd}
\pdfglyphtounicode{parenleftbig}{0028}
\pdfglyphtounicode{parenrightbig}{0029}
\hypersetup{pdftitle={Schnell fiber spaces and good canonical models},pdfauthor={OpenAI}}
\pdfinfoomitdate=1
\pdftrailerid{}
\newtheorem{theorem}{Theorem}[section]
\newtheorem{lemma}[theorem]{Lemma}
\newtheorem{proposition}[theorem]{Proposition}
\newtheorem{corollary}[theorem]{Corollary}
\newtheorem{assumption}[theorem]{Assumption}
\theoremstyle{definition}
\newtheorem{definition}[theorem]{Definition}
\newtheorem{remark}[theorem]{Remark}
\newcommand{\C}{\mathbb C}
\newcommand{\Q}{\mathbb Q}
\newcommand{\R}{\mathbb R}
\newcommand{\Z}{\mathbb Z}
\newcommand{\N}{\mathbb N}
\newcommand{\cO}{\mathcal O}
\newcommand{\PE}{\overline{\operatorname{Eff}}}
\DeclareMathOperator{\Supp}{Supp}
\DeclareMathOperator{\codim}{codim}
\DeclareMathOperator{\rk}{rk}
\DeclareMathOperator{\Pic}{Pic}
\numberwithin{equation}{section}
\setlist[enumerate]{itemsep=3pt,topsep=5pt}
\setlist[itemize]{itemsep=3pt,topsep=5pt}
\title{Schnell fiber spaces and good canonical models}
\author{OpenAI}
\date{September 24, 2026}
\begin{document}
\maketitle
\begin{abstract}
We prove the Campana--Peternell inequality \(\kappa(X)\geq\kappa(D)\) for a
smooth connected projective complex variety \(X\) and an effective Cartier
divisor \(D\) whenever \(m_0K_X-D\) is pseudo-effective for some positive
integer \(m_0\). For an algebraic fiber space \(f\colon X\to Y\) between
smooth connected projective complex varieties, the hypothesis that
\(m_0K_X-f^*H\) is pseudo-effective with \(H\) ample Cartier gives
\(\kappa(X)=\kappa(F)+\dim Y\) for a very general smooth fiber \(F\), as well
as nonzero sections of \(mK_X-f^*H\) for all sufficiently large divisible
\(m\).
\end{abstract}
\section{Introduction}
\label{sec:introduction}

The pluricanonical systems of an algebraic fiber space reflect both
the geometry of its fibers and the positivity available from its base.
When the geometric generic fiber has Kodaira dimension zero, the
easy-addition bound says that the total space has Kodaira dimension
at most the dimension of the base. The zero-Kodaira case of Schnell's
question asks whether a numerical comparison with an ample divisor on
the base forces equality.
Here an \emph{algebraic fiber space} means a surjective morphism with
connected fibers between projective varieties.

We approach this question through a good canonical model of the total
space. The geometric argument below shows what such a model supplies;
the smooth canonical good-model theorem of
\cite[Corollary~11.2]{LogAbundance2026} then provides the model.
Its nonvanishing consequence also permits Schnell's reduction to pass
from zero-Kodaira fibers to the Campana--Peternell inequality and the
general fiber-space conclusion stated below.

The formulation in \cite[Conjecture~1.2]{Kim2025} compares
a positive multiple of $K_X$ with the pullback of an ample Cartier divisor on the base.
The difficulty is that pseudo-effectivity concerns a numerical divisor
class, whereas Kodaira dimension measures actual sections. The
following theorem gives the positive conclusion in this formulation;
its assertion has also been proved by Zou, as discussed below.

\begin{theorem}[Schnell's zero-Kodaira fiber-space conclusion]
\label{thm:schnell}
Let $f\colon X\to Y$ be a surjective morphism with connected fibers
between smooth connected projective complex varieties. Let $F$ be its
smooth geometric generic fiber and assume $\kappa(F)=0$. Suppose that
there are an ample Cartier divisor $H$ on $Y$ and a positive integer
$m_0$ such that $m_0K_X-f^*H$ is pseudo-effective. Then
\[
 \kappa(X)=\dim Y.
\]
\end{theorem}

The ample comparison is essential. For example, let $F$ be smooth,
connected and projective with torsion canonical bundle. The projection
$F\times\mathbb P^1\to\mathbb P^1$ has fibers of Kodaira dimension
zero, but its total space has Kodaira dimension $-\infty$.
The class $m_0K_{F\times\mathbb P^1}-f^*H$ has negative degree on
the moving curves $\{x\}\times\mathbb P^1$ for every ample $H$;
it therefore cannot be pseudo-effective: for any fixed effective
divisor, a general member of this family is not contained in its
support and has nonnegative intersection with it. In contrast, for a
smooth connected projective $Y$ with $K_Y$ ample, the projection
$F\times Y\to Y$ satisfies the numerical
hypothesis with $H=K_Y$ for every $m_0\geq1$, since the compared
class is numerically $(m_0-1)f^*K_Y$. Its Kodaira dimension is
$\dim Y$.

The canonical good-model theorem used below also supplies canonical
nonvanishing in every dimension. With this input, Schnell's reduction
\cite[Sections~4--10]{Schnell2022} promotes
Theorem~\ref{thm:schnell} to the Kodaira-dimension form of the
Campana--Peternell conjecture.

\begin{corollary}[Campana--Peternell]
\label{cor:campana-peternell}
Let $X$ be a smooth connected projective complex variety, let $D$ be
an effective Cartier divisor on $X$, and let $m_0$ be a positive
integer. If $m_0K_X-D$ is pseudo-effective, then
\[
 \kappa(X)\geq\kappa(D).
\]
\end{corollary}

The same reduction gives the general fiber-space conclusion, including
the existence of sections after subtracting the original ample
pullback.

\begin{corollary}[Schnell's general fiber-space conclusion]
\label{cor:schnell-general}
Let $f\colon X\to Y$ be a surjective morphism with connected fibers
between smooth connected projective complex varieties, and let $F$ be
a very general smooth fiber. Suppose that $H$ is an ample Cartier
divisor on $Y$ and $m_0$ is a positive integer such that
$m_0K_X-f^*H$ is pseudo-effective. Then
\[
 \kappa(X)=\kappa(F)+\dim Y.
\]
Moreover, there are positive integers $r$ and $\ell_0$ such that
\[
 H^0\bigl(X,\cO_X(\ell rK_X-f^*H)\bigr)\ne0
 \qquad\text{for every integer }\ell\geq\ell_0.
\]
\end{corollary}

The last assertion is effectivity for all sufficiently large and
divisible pluricanonical degrees, with no uniform choice of $r$ or
$\ell_0$ asserted. The numerical hypothesis itself implies
$\kappa(F)\geq0$ by restriction and canonical nonvanishing; it is not
replaced by that weaker fiber condition. Corollary~\ref{cor:campana-peternell}
also applies to an effective rational divisor after clearing
denominators, as explained in Section~\ref{sec:general-consequences}.

\subsection{The problem and previous approaches}

Campana and Peternell formulated the comparison
$NK_X=A+B$, with $N$ a positive integer, $A$ effective and $B$
pseudo-effective, in their
study of positivity of the cotangent bundle
\cite[Conjecture~2.4]{CP2011}. The proposed inequality
$\kappa(X)\geq\kappa(A)$ includes canonical nonvanishing when
$A=0$: pseudo-effective $K_X$ should have a nonzero pluricanonical
section. For general $A$, the conclusion also requires the canonical
systems to have at least the image dimension supplied by $A$.
When $\kappa(X)\geq0$, Campana and Peternell reduce this comparison
to the general fibers of the Iitaka fibration of $X$, which have
Kodaira dimension zero \cite[Proposition~2.6]{CP2011}. In that
zero-Kodaira setting, a good minimal model has torsion canonical
class; their pushforward argument then forces $\kappa(A)=0$
\cite[Proposition~2.7]{CP2011}. This already exhibits how a good
model can convert numerical information into a statement about
sections.

Schnell developed this circle of questions in his study of singular
metrics, nonvanishing and the Campana--Peternell conjecture
\cite[Sections~4--10]{Schnell2022}. His Conjecture~10.1 allows fibers
of nonnegative Kodaira dimension and asks for nonvanishing after
subtracting an ample pullback. Under canonical nonvanishing, his
Section~9 reduces the conjecture to fibers of Kodaira dimension zero;
Lemma~7.1 supplies an effective pluricanonical divisor after a positive
base twist. If the fiber has positive Kodaira dimension, the Iitaka
fibration of that divisor has a strictly larger base, while retaining
the numerical comparison. Iteration reaches the zero-Kodaira case.
Section~8 connects the resulting equality to effectivity after
subtracting the original ample pullback. This uses the Iitaka-addition
criterion of Fujita and Mori, in the form recorded in
\cite[Lemma~4.6]{PopaSchnell2014}. He proves a special case when the
canonical divisor of the base is pseudo-effective
\cite[Theorem~12.1]{Schnell2022}.

Kim applies the canonical bundle formula to the same fiber-space
problem. His Theorem~1.3 proves the conclusion when
$K_Y+(1-\varepsilon)B_Y$ is pseudo-effective for some
$\varepsilon>0$, where $B_Y$ is the discriminant on the birational
setup of that theorem, or when the canonical class of the general
fiber is rigid \cite{Kim2025}. Here rigidity means uniqueness of
the closed positive $(1,1)$-current representing that class.
Kim also recalls a known algebraic proof under the existence of a
good minimal model of the general fiber
\cite[Section~1.3]{Kim2025}. Zou obtains the same
conclusion as Theorem~\ref{thm:schnell}, without these additional
hypotheses, by a canonical-bundle-formula argument
\cite[Theorems~1.3 and~5.1]{Zou2025}.

We give a direct mixed-intersection proof of
Theorem~\ref{thm:schnell} from a good minimal model of the
\emph{total space}. We separate that geometric argument from the
model-existence theorem that it uses. The latter is the smooth
canonical good-model theorem of
\cite[Corollary~11.2]{LogAbundance2026}, quoted in
Theorem~\ref{thm:canonical-good-models}. Its nonvanishing consequence
and Schnell's published reduction then yield
Corollaries~\ref{cor:campana-peternell} and~\ref{cor:schnell-general}.

\subsection{How the numerical condition produces the lower bound}

A good klt minimal model of $X$ consists of a normal projective
$\Q$-factorial klt variety $V$ with semiample canonical divisor and a
$K_X$-negative birational contraction $X\dashrightarrow V$.
For compatible canonical divisors, the comparison on a smooth common
resolution $X\xleftarrow{p}W\xrightarrow{q}V$ has the form
\begin{equation}\label{eq:intro-comparison}
 p^*K_X=q^*K_V+E,\qquad E\geq0,\qquad q_*E=0.
\end{equation}
The last condition means that every component of $E$ is exceptional
over $V$. The model therefore provides sections and also specifies
the error in transferring them to $X$.

Choose a globally generated Cartier multiple $rK_V$. The comparison
makes $rE$ an effective Cartier divisor. Pulling back sections of
$rK_V$ and multiplying by the canonical section of $rE$ gives a
pluricanonical subsystem on $X$ with the same image dimension.
It remains to force that dimension to be at least $\dim Y$.

If the image were smaller, Lemma~\ref{lem:mixed-positivity} would
provide a mixed complete-intersection curve class $C$ on $W$ with
three properties: it pairs nonnegatively with every pseudo-effective
divisor, it annihilates both $q^*K_V$ and $E$, and it pairs positively
with $(fp)^*H$. Its divisor factors are pulled back from $V$, so
the exceptional vanishing follows from the Cartier projection formula.
The strict positivity uses a cotangent direction from the map to $Y$
beyond those supplied by the canonical morphism. Thus
\[
 (m_0p^*K_X-(fp)^*H)\cdot C=-(fp)^*H\cdot C<0,
\]
contrary to the pulled-back numerical hypothesis. Section~\ref{sec:geometry}
constructs the class and proves each of these properties.

This argument uses the classical intersection theory of Cartier
divisors and proper pushforward \cite[Chapter~2]{Fulton1998}.
In particular, the maps to $Y$ and to the canonical image need no factorization
relation. Section~\ref{sec:geometry} gives the complete geometric
proof, including the upper bound and the transfer of sections.
Section~\ref{sec:good-model-interface} states the exact good-model
input, verifies its hypotheses and completes
Theorem~\ref{thm:schnell}. Section~\ref{sec:general-consequences}
explains the nonvanishing and fiber interfaces, then proves the
Campana--Peternell and general Schnell conclusions.

\subsection{Conventions}

All varieties are integral and projective over $\C$ unless specified
otherwise. Smoothness of an algebraic fiber space refers here to its
source and target; special fibers of the morphism may be singular.
Canonical divisors on birational models are chosen compatibly.
For a rational Cartier divisor $L$, the Iitaka dimension $\kappa(L)$
is the maximum dimension of the rational images of its nonempty
complete systems $|mL|$, with $m>0$ clearing the Cartier index.
It is $-\infty$ if all these systems are empty. We write
$\kappa(X)=\kappa(K_X)$.

We write $\PE(X)$ for the closure of the cone of effective real
divisor classes in $N^1(X)_\R$. A divisor is pseudo-effective when
its numerical class belongs to this cone. A rational Cartier divisor
is semiample when a positive Cartier multiple is globally generated.
The section constructions below use this actual line-bundle property.

\section{The numerical argument on a good minimal model}
\label{sec:geometry}

We prove the fiber-space conclusion assuming that the total space has
a good klt minimal model. The upper bound comes from the geometric
generic fiber. For the lower bound, we construct a mixed intersection
that detects the ample class from the base and kills both terms in
the canonical comparison.

\begin{theorem}[The good-model case]\label{thm:good-model-case}
Let $f\colon X\to Y$ be a surjective morphism with connected fibers
between smooth connected projective complex varieties. Let $F$ be its
smooth geometric generic fiber and suppose $\kappa(F)=0$. Assume
that $X$ has a projective good klt minimal model. If $H$ is an ample
Cartier divisor on $Y$ and $m_0$ is a positive integer such that
\begin{equation}\label{eq:geometric-numerical-hypothesis}
 m_0K_X-f^*H\in\PE(X),
\end{equation}
then $\kappa(X)=\dim Y$.
\end{theorem}

We use the good-model convention and effective comparison
\eqref{eq:intro-comparison} stated in the introduction. This is the
$K_X$-negative convention of \cite[Definitions~3.6.1 and~3.6.7]{BCHM2010}.

\subsection{The upper bound from the generic fiber}

We first account for the use of $\kappa(F)=0$. It bounds every
pluricanonical image of $X$, before any model is chosen.

\begin{lemma}[The geometric generic fiber bound]\label{lem:easy-addition}
Let $f\colon X\to Y$ be a surjective morphism with connected fibers
between smooth connected projective complex varieties. If its smooth
geometric generic fiber $F$ satisfies $\kappa(F)=0$, then
$\kappa(X)\leq\dim Y$.
\end{lemma}

\begin{proof}
Put $K=\C(Y)$ and $\eta=\operatorname{Spec}K$. Generic smoothness
and Stein factorization give a smooth geometrically integral generic
fiber $X_\eta$ \cite[III, Sections~10--11]{Hartshorne1977}.
The cotangent sequence along this fiber gives
\begin{equation}\label{eq:canonical-generic-restriction}
 \omega_X|_{X_\eta}
 \simeq\omega_{X_\eta/K}\otimes_K\ell,
 \qquad \ell=\omega_Y|_\eta.
\end{equation}
Here $\ell$ is a one-dimensional $K$-vector space. The factor from
the base therefore cancels in ratios of pluricanonical sections.

Fix $m>0$ such that $H^0(X,mK_X)\ne0$, and choose a basis
$s_0,\ldots,s_N$ with $s_0\ne0$. A nonzero section is nonzero at
the generic point of $X$, so it stays nonzero on $X_\eta$ and after
extension to the geometric generic fiber. By
\eqref{eq:canonical-generic-restriction}, the restricted system on
$F$ is a nonzero subsystem of $|mK_F|$. Its image has dimension zero,
since $\kappa(F)=0$.

Let $\phi_m$ be the rational map of the complete system $|mK_X|$,
and let $T$ be the closure of the image of
$(f,\phi_m)\colon X\dashrightarrow Y\times\mathbb P^N$.
Its function field is
\[
 L=K(s_1/s_0,\ldots,s_N/s_0)\subseteq\C(X).
\]
The dimension of its generic image over $Y$ is
$\operatorname{trdeg}_K L$. This dimension is unchanged by algebraic
extension of $K$. Geometric integrality of $X_\eta$ identifies that
extended image with the image of the same section ratios on $F$.
Thus $\operatorname{trdeg}_K L=0$, and $\dim T=\dim Y$.
Projection to $\mathbb P^N$ gives $\dim\phi_m(X)\leq\dim Y$.
Taking the maximum over all nonempty pluricanonical systems proves
the lemma. If every system is empty, the inequality is immediate.
\end{proof}

\subsection{A curve class that detects the base}

The lower bound requires a different ingredient. The following lemma
constructs a numerical test from a globally generated divisor on $V$.
All its factors come from $V$, which is why exceptional errors vanish.
Its strict positivity uses the independent morphism to $Y$.

\begin{lemma}[A mixed intersection test]\label{lem:mixed-positivity}
Let $q\colon W\to V$ be a birational morphism from a smooth
projective complex $d$-fold to a normal projective variety.
Let $D$ be a globally generated Cartier divisor on $V$, and let
$g\colon V\to Z\subseteq\mathbb P^N$ be the morphism of $|D|$
onto its image. Write $k=\dim Z$, and let $A$ be a very ample
Cartier divisor on $V$. Suppose $h\colon W\to Y$ is a surjective
morphism to a smooth projective variety with $y=\dim Y>k$.
For every ample Cartier divisor $H$ on $Y$, the numerical one-cycle
class
\[
 C=(q^*D)^k(q^*A)^{d-k-1}
\]
satisfies the following properties:
\begin{enumerate}
 \item $B\cdot C\geq0$ for every pseudo-effective real divisor
 class $B$ on $W$;
 \item $E\cdot C=0$ for every $q$-exceptional rational divisor $E$;
 \item $q^*D\cdot C=0$ and $h^*H\cdot C>0$.
\end{enumerate}
\end{lemma}

\begin{proof}
The inequality $k<y\leq d$ makes all exponents nonnegative; a
product with no factors has its usual meaning. We prove first that
the test is nonnegative, then that it has the two required vanishing
properties, and finally that it detects the ample divisor from $Y$.

\smallskip\noindent
\emph{Nonnegativity.}
The divisors $q^*D$ and $q^*A$ are globally generated. If $B$ is
effective, choose their members successively so that no member
contains an irreducible component of the preceding intersection
on $B$. Global generation permits all the finitely many required
avoidances. The resulting Cartier intersections form an effective
zero-cycle, possibly empty. Hence $B\cdot C\geq0$
\cite[Chapter~2]{Fulton1998}. Extend by linearity to effective real
divisors and by continuity on $N^1(W)_\R$ to its closed effective
cone. This proves (1).

\smallskip\noindent
\emph{Vanishing on exceptional and semiample classes.}
For a $q$-exceptional divisor the Cartier projection formula gives
\begin{equation}\label{eq:exceptional-annihilation}
 E\cdot(q^*D)^k(q^*A)^{d-k-1}
   =q_*E\cdot D^kA^{d-k-1}=0;
\end{equation}
see \cite[Proposition~2.3(c)]{Fulton1998}.
Indeed, every component of $E$ maps to codimension at least two,
so $q_*E=0$ as a divisor cycle. This reasoning is valid on the
possibly singular variety $V$.
Also $D=g^*\cO_Z(1)$, and $k+1$ general hyperplanes miss the
$k$-dimensional image $Z$. Consequently $(q^*D)^{k+1}=0$.
For $k=0$, the morphism $g$ is constant and $\cO_V(D)$ is trivial,
which gives the same conclusion. These observations prove (2) and
the first assertion of (3).

\smallskip\noindent
\emph{Strict positivity.}
Choose $b>0$ such that $bH$ is very ample. Use this divisor to
embed $Y$, and use $A$ to embed $V$.
There is a point $w\in W$ at which $q$ is an isomorphism onto a
smooth open subset of $V$ and
\[
 \operatorname{rank}d(gq)_w=k,
 \qquad\operatorname{rank}dh_w=y,
 \qquad\operatorname{rank}dq_w=d.
\]
Each condition holds on a dense open subset. For the rank statements,
the generic differential rank in characteristic zero equals the
dimension of the image, by separability of the function-field
extension. We choose $w$ in the intersection of these open subsets.

The differentials of pulled-back hyperplanes through $gq(w)$ span
a $k$-dimensional subspace $U\subseteq T_w^*W$. Choose $k$ of them
whose differentials form a basis of $U$. The hyperplanes through
$h(w)$ supply a $y$-dimensional subspace. Since $y>k$, one such
hyperplane pulls back to a divisor with differential outside $U$.
Finally, hyperplanes from the embedding by $A$ supply all cotangent
directions at $w$, because $q$ is a local isomorphism there. Choose
$d-k-1$ of them completing the chosen differentials to a basis of
$T_w^*W$. These $d$ divisors meet transversely at $w$.

To use this local intersection in the global intersection number,
perturb the hyperplanes slightly in their complex parameter spaces.
The transverse point persists for every sufficiently small
perturbation, by the implicit function theorem. Tuples whose
pullbacks meet properly at every successive step form a nonempty
Zariski open subset of the product of the hyperplane parameter
spaces. Nonemptiness follows by successively avoiding the finitely
many components already present, using global generation; openness
follows from upper semicontinuity of fiber dimension for the
projective universal intersections. This parameter space is
irreducible, so the open subset is dense also in the complex
analytic topology. It therefore meets the perturbation neighborhood.

The resulting proper global intersection is an effective zero-cycle
that contains a transverse point of multiplicity one. It follows that
\[
 h^*(bH)\cdot(q^*D)^k(q^*A)^{d-k-1}>0.
\]
Division by $b$ proves (3). Only one cotangent direction from $h$
outside $U$ was needed; the two morphisms $gq$ and $h$ need not
factor through each other.
\end{proof}

\subsection{Sections and the numerical contradiction}

We now combine the two lemmas. The good model supplies actual
sections on $X$. The mixed test forces their image to have at least
the dimension of $Y$.

\begin{proof}[Proof of Theorem~\ref{thm:good-model-case}]
If $Y$ is a point, then $F=X$ and $\kappa(F)=0$ is the required
conclusion. Hence assume $d=\dim X\geq y=\dim Y>0$.

Let $V$ be a good klt minimal model of $X$. Resolve the closure of
the graph of $X\dashrightarrow V$ projectively
\cite{Hironaka1964}, obtaining a smooth projective common resolution
$X\xleftarrow{p}W\xrightarrow{q}V$. By the good-model comparison,
compatible canonical divisors satisfy
\begin{equation}\label{eq:canonical-comparison}
 p^*K_X=q^*K_V+E,\qquad E\geq0,\qquad q_*E=0.
\end{equation}
Pulling back to a further common resolution preserves the effectivity
and target-exceptionality of this error.

Choose $r>0$ such that $D=rK_V$ is Cartier and globally generated,
and let $g\colon V\to Z\subseteq\mathbb P^N$ be its morphism onto
its image. Put $k=\dim Z$. Since $K_X$ is Cartier,
\eqref{eq:canonical-comparison} shows that
$rE=p^*(rK_X)-q^*D$ is an integral Cartier divisor on $W$.
The two maps whose dimensions we will compare appear in
Figure~\ref{fig:common-resolution}.

\begin{figure}[ht]
\centering
\begin{tikzcd}[column sep=large,row sep=large]
 & W \arrow[dl,"p"'] \arrow[dr,"q"] & \\
 X \arrow[d,"f"'] && V \arrow[d,"g"] \\
 Y && Z
\end{tikzcd}
\caption{The common resolution carries the two morphisms $h=fp$ to
$Y$ and $gq$ to $Z$.
The divisor factors defining the mixed class are pulled back from
$V$, so the class annihilates the
$q$-exceptional error. Its pairing with $h^*H$ detects the ample
class from the separate base $Y$.}
\label{fig:common-resolution}
\end{figure}

Pull back sections of $D$ by $q$ and multiply by the canonical
section of the effective Cartier divisor $rE$. We obtain an injection
\begin{equation}\label{eq:section-transfer}
 H^0(V,D)\hookrightarrow H^0(W,rp^*K_X)=H^0(X,rK_X).
\end{equation}
The equality follows from $p_*\cO_W=\cO_X$ and the sheaf projection
formula for the proper birational map to the normal variety $X$
\cite[III, Corollary~11.4 and its proof; II, Exercise~5.1(d)]{Hartshorne1977}.
Multiplication by the section of $rE$ leaves section ratios unchanged
on the complement of its support. Thus the subsystem in
\eqref{eq:section-transfer} has image dimension $k$, and
\begin{equation}\label{eq:canonical-image-lower-bound}
 \kappa(X)\geq k.
\end{equation}

Suppose, for a contradiction, that $k<y$. Choose a very ample
Cartier divisor $A$ on $V$ and apply
Lemma~\ref{lem:mixed-positivity} with
\[
 h=fp,\qquad C=(q^*D)^k(q^*A)^{d-k-1}.
\]
The canonical comparison and the vanishing assertions of the lemma
give
\begin{equation}\label{eq:canonical-test-zero}
 p^*K_X\cdot C
 =\frac1r q^*D\cdot C+E\cdot C=0.
\end{equation}
Pullback by $p$ preserves pseudo-effectivity: on smooth varieties
it takes effective real divisors to effective real divisors and
induces a continuous linear map on numerical divisor spaces.
Therefore the pullback of
\eqref{eq:geometric-numerical-hypothesis}, paired with $C$, yields
\[
 0\leq(m_0p^*K_X-h^*H)\cdot C=-h^*H\cdot C<0,
\]
a contradiction. Hence $k\geq y$. Equation
\eqref{eq:canonical-image-lower-bound} and
Lemma~\ref{lem:easy-addition} now give
$y\leq k\leq\kappa(X)\leq y$, proving the theorem.
\end{proof}

The proof includes $k=0$, $d-k-1=0$ and relative dimension zero.
Effectivity of $E$ supplies the section injection, while
target-exceptionality makes its mixed pairing vanish. These are
separate uses of the canonical comparison, and both are needed.

\section{The canonical good model and the main theorem}
\label{sec:good-model-interface}

The geometric argument has used the existence of a good model of
$X$ through its sections and its effective exceptional comparison.
We now supply exactly that model. The following theorem is
\cite[Corollary~11.2]{LogAbundance2026}; its proof belongs to that
companion article.

\begin{theorem}[Smooth canonical good models]
\label{thm:canonical-good-models}
Let $T$ be a smooth connected projective complex variety with
pseudo-effective $K_T$. Then there exist a normal projective
$\Q$-factorial klt variety $V$ and a $K_T$-negative birational
contraction $T\dashrightarrow V$ such that $K_V$ is $\Q$-Cartier
and semiample. On a smooth projective common resolution
$T\xleftarrow{p}W\xrightarrow{q}V$, compatible canonical divisors
satisfy
\[
 p^*K_T=q^*K_V+E,\qquad E\geq0,\qquad q_*E=0.
\]
\end{theorem}

Thus the input provides both the actual globally generated multiple
used for sections and the exceptional equality used for intersections.

\begin{proof}[Proof of Theorem~\ref{thm:schnell}]
If $Y$ is a point, then its geometric generic fiber is $X$, and
the conclusion is the hypothesis $\kappa(F)=0$.
For a positive-dimensional base, a positive multiple of the ample
divisor $H$ has an effective representative. Its pullback shows that $f^*H$ is
pseudo-effective. Adding it to the original numerical hypothesis gives
\[
 K_X=\frac1{m_0}\bigl((m_0K_X-f^*H)+f^*H\bigr)\in\PE(X).
\]
The variety $X$ is smooth, connected, projective and complex, so
Theorem~\ref{thm:canonical-good-models} applies. Its output is the
projective good klt minimal model required by
Theorem~\ref{thm:good-model-case}. Applying that theorem to the
original $f$, $H$ and $m_0$ gives $\kappa(X)=\dim Y$.
\end{proof}

Throughout the proof the assumed comparison is numerical. The
pluricanonical sections are obtained from the semiample model through
\eqref{eq:section-transfer}; they are not assumed as a reformulation
of pseudo-effectivity.

\section{Campana--Peternell and general Schnell fiber spaces}
\label{sec:general-consequences}

The good-model input supplies canonical nonvanishing as well as the
model used in the mixed-intersection argument. These two outputs fit
the reduction of \cite[Sections~4--10]{Schnell2022}.

\subsection{Nonvanishing and very general fibers}

Let $T$ be a smooth connected projective complex variety with
pseudo-effective $K_T$, and take the model $V$ and comparison divisor
$E$ supplied by Theorem~\ref{thm:canonical-good-models}.
Choose a positive multiple $rK_V$ that is Cartier and globally
generated. The comparison then makes $rE$ an effective integral
Cartier divisor. The section transfer
in~\eqref{eq:section-transfer} gives
\[
 0\ne H^0(V,\cO_V(rK_V))
 \lhook\joinrel\longrightarrow H^0(T,\cO_T(rK_T)).
\]
Consequently,
\begin{equation}\label{eq:canonical-nonvanishing}
 K_T\in\PE(T)\quad\Longrightarrow\quad \kappa(T)\geq0.
\end{equation}
This is the canonical nonvanishing input used by Schnell.

We will use smooth closed complex fibers when applying
\eqref{eq:canonical-nonvanishing}. They have the same Kodaira
dimension as the geometric generic fiber when chosen very generally.
Indeed, let $f\colon X\to Y$ be a surjective morphism with connected
fibers between smooth connected projective complex varieties, and
write $X_{\bar\eta}$ for its geometric generic fiber. Over a nonempty
open subset, $f$ is smooth; there
$\omega_{X/Y}=\omega_X\otimes f^*\omega_Y^{-1}$ restricts to the
canonical bundle of each fiber. For each positive integer $m$, after
shrinking this open subset, formation of
$f_*(\omega_{X/Y}^{\otimes m})$ commutes with base change.
Outside the union of the resulting countably many proper
closed subsets, a smooth fiber $F$ therefore satisfies
\[
 h^0(F,\omega_F^{\otimes m})
 =
 h^0(X_{\bar\eta},\omega_{X_{\bar\eta}}^{\otimes m})
 \qquad\text{for every }m>0.
\]
Here the dimension on the right is over the algebraic closure of
$\C(Y)$. Thus the plurigenus sequences, and hence the Kodaira
dimensions, agree. Such complex fibers exist because $\C$ is
uncountable. In particular, the zero-Kodaira fiber condition in the
reduction below is precisely the geometric generic fiber condition
of Theorem~\ref{thm:schnell}.

We also use the restriction observation in
\cite[Section~7]{Schnell2022}. A pseudo-effective real Cartier
divisor $L$ on $X$ restricts to a pseudo-effective divisor on a very
general fiber. To see this, choose effective real divisors whose
numerical classes converge to $[L]$. After excluding a countable
union of proper closed subsets of $Y$, a fiber is contained in none
of their supports. Their restrictions are effective, and their
numerical classes converge to the class of $L|_F$. Applying this to
$L=m_0K_X-f^*H$ gives
\[
 K_F\in\PE(F),\qquad \kappa(F)\geq0,
\]
where the second assertion follows from
\eqref{eq:canonical-nonvanishing}. The very general fiber can be
chosen to satisfy this restriction property and the plurigenus
comparison simultaneously.

\subsection{The Campana--Peternell inequality}

\begin{proof}[Proof of Corollary~\ref{cor:campana-peternell}]
Since $D$ is effective, $\kappa(D)\geq0$. Also
\[
 K_X=\frac1{m_0}\bigl((m_0K_X-D)+D\bigr)\in\PE(X).
\]
If $\kappa(D)=0$, canonical nonvanishing
\eqref{eq:canonical-nonvanishing} proves the assertion. We may
therefore assume $\kappa(D)>0$.

Schnell's reduction in \cite[Sections~4--6]{Schnell2022} first
resolves a system $|nD|$ whose image has dimension $\kappa(D)$,
then takes its Stein factorization and resolves the base. It produces
an algebraic fiber space
$f_0\colon X_0\to Y_0$ between smooth connected projective complex
varieties, with $X_0$ birational to $X$, an ample Cartier divisor
$H_0$ on $Y_0$, and a positive integer $a_0$ such that
\begin{equation}\label{eq:initial-cp-reduction}
 \dim Y_0=\kappa(D),\qquad
 a_0K_{X_0}-f_0^*H_0\in\PE(X_0).
\end{equation}
The effective exceptional terms from resolving $X$ preserve
pseudo-effectivity; after resolving the base, Schnell replaces the
big and nef pullback of the original ample divisor by a suitable
ample divisor. In particular, the output has the precise ample
Cartier hypothesis of Theorem~\ref{thm:schnell}.

We recall the finite iteration in
\cite[Sections~7 and~9]{Schnell2022}. Suppose that
$f_i\colon X_i\to Y_i$ is such a fiber space, with ample Cartier
$H_i$ and
$a_iK_{X_i}-f_i^*H_i\in\PE(X_i)$ for a positive integer $a_i$.
For a very general smooth fiber $F_i$, the preceding restriction and
nonvanishing argument gives $\kappa(F_i)\geq0$. If
$\kappa(F_i)>0$, the construction in
\cite[Lemma~7.1]{Schnell2022} chooses positive integers $r_i,b_i$
and an effective Cartier divisor $L_i$ such that
\[
 L_i\sim r_iK_{X_i}+f_i^*(b_iH_i),\qquad
 \kappa(L_i)=\kappa(F_i)+\dim Y_i.
\]
The resulting comparison is
\[
 (a_ib_i+r_i)K_{X_i}-L_i
 \sim b_i(a_iK_{X_i}-f_i^*H_i).
\]
Its right side is pseudo-effective, so the same is true of its left
side, since linear equivalence preserves numerical classes.
Applying the divisor-to-fiber-space reduction of Sections~4--6 of
Schnell's paper to $L_i$ gives another algebraic fiber space
$f_{i+1}\colon X_{i+1}\to Y_{i+1}$ of the same smooth projective
complex type, with $X_{i+1}$ birational to $X_i$, an ample Cartier
divisor $H_{i+1}$, and a positive integer $a_{i+1}$ such that
\[
 \dim Y_{i+1}=\kappa(F_i)+\dim Y_i,\qquad
 a_{i+1}K_{X_{i+1}}-f_{i+1}^*H_{i+1}\in\PE(X_{i+1}).
\]

As long as $\kappa(F_i)>0$, the integer $\dim Y_i$ strictly
increases. It is bounded by $\dim X$, so this process reaches an
index $j$ with $\kappa(F_j)=0$. The geometric generic fiber has
Kodaira dimension zero by the comparison above.
Theorem~\ref{thm:schnell} applies at this last stage. Birational
invariance of Kodaira dimension and
\eqref{eq:initial-cp-reduction} give
\[
 \kappa(X)=\kappa(X_j)=\dim Y_j
 \geq\dim Y_0=\kappa(D).
\]
\end{proof}

If $D$ is an effective rational divisor instead, choose a positive
integer $q$ such that $qD$ is an integral Cartier divisor.
Multiplying the numerical hypothesis by $q$ gives
$qm_0K_X-qD\in\PE(X)$. Corollary~\ref{cor:campana-peternell}
applied to $qD$ gives the same conclusion because
$\kappa(qD)=\kappa(D)$. This extension still requires an effective
divisor; it makes no assertion for an arbitrary pseudo-effective
divisor in its place.

\subsection{Equality and eventual effectivity for general fibers}

\begin{proof}[Proof of Corollary~\ref{cor:schnell-general}]
If $Y$ is a point, then $F=X$ and $H$ is linearly equivalent to zero.
The equality is tautological, and the numerical hypothesis gives
pseudo-effective $K_X$. A nonzero section supplied by
\eqref{eq:canonical-nonvanishing}, together with its powers, gives
the asserted sections.

Assume $\dim Y>0$. The restriction argument above gives
$\kappa(F)\geq0$. As in \cite[Lemma~7.1]{Schnell2022}, choose
positive integers $a,b$ and an effective Cartier divisor $L$ with
\[
 L\sim aK_X+f^*(bH),\qquad \kappa(L)=\kappa(F)+\dim Y.
\]
The linear equivalence
\[
 (m_0b+a)K_X-L\sim b(m_0K_X-f^*H)
\]
shows that its left side is pseudo-effective.
Corollary~\ref{cor:campana-peternell}, applied to the effective
Cartier divisor $L$, gives
\[
 \kappa(X)\geq\kappa(L)=\kappa(F)+\dim Y.
\]
The reverse inequality is the easy-addition inequality used in
\cite[Lemma~7.1]{Schnell2022}. This proves the equality.

To obtain sections after subtracting the original $f^*H$, we use
the Fujita--Mori criterion recalled in
\cite[Lemma~4.6]{PopaSchnell2014}; this is the criterion behind
\cite[Section~8]{Schnell2022}. Since $\kappa(F)\geq0$, the equality
just proved implies that there are a big Cartier divisor $B$ on $Y$
and a positive integer $u$ with a nonzero section
\[
 \sigma\in H^0\bigl(X,\cO_X(uK_X-f^*B)\bigr).
\]
Bigness of $B$ gives a positive integer $v$ and a nonzero section
$\tau\in H^0(Y,\cO_Y(vB-H))$. Thus, with $r=uv$,
\[
 s=\sigma^v f^*\tau
 \in H^0\bigl(X,\cO_X(rK_X-f^*H)\bigr)
\]
is nonzero. This step replaces the big divisor supplied by the
criterion with the ample divisor in the original hypothesis.

Finally, ampleness of $H$ supplies a nonzero section
$\tau_\ell\in H^0(Y,\cO_Y((\ell-1)H))$ for every sufficiently
large integer $\ell$. The products
\[
 s^\ell f^*\tau_\ell
 \in H^0\bigl(X,\cO_X(\ell rK_X-f^*H)\bigr)
\]
are nonzero. Choosing $\ell_0$ beyond this ampleness threshold proves
the asserted conclusion for every integer $\ell\geq\ell_0$.
\end{proof}

\bibliographystyle{plain}
\bibliography{references}
\end{document}
